\section{Introduction}

The integration of Large Language Models (LLMs) into the scientific workflow represents a paradigm shift in theoretical physics. This graduate seminar explores how to leverage advanced Artificial Intelligence not merely as a basic computational utility, but as an active collaborator in the research process. Designed for graduate students seeking to modernize their scientific toolkit, this course focuses on three core pillars: the rapid assimilation of novel physical concepts, advanced theoretical problem-solving, and the critical evaluation of AI-generated insights.

Rather than a traditional lecture format, this seminar is highly interactive and project-driven. Students will tackle a series of complex, multi-stage theoretical problems spanning different subfields of physics. Working in collaborative teams, participants will use LLMs to navigate unfamiliar literature, build intuition for new topics, formulate mathematical models, and write simulation code.

Crucially, the course places a strong emphasis on the critical use of AI. Students will learn the limitations of current models – specifically how to rigorously verify AI outputs, identify physical inconsistencies or "hallucinations", and design effective prompts for highly technical, mathematical subjects.


\subsection{Course Information}

\begin{itemize}
\item \textbf{Format:} Weekly 2-hour discussion and presentation sessions, supplemented by collaborative, hackathon-style teamwork outside of class. Submission of a final report (individually).
\item \textbf{Prerequisites:} A solid graduate-level foundation in core physics. No extensive prior knowledge of the specific problem domains is required; the course will teach you how to use AI to rapidly build the necessary theoretical scaffolding.
\item \textbf{Learning Outcomes:} By the end of the term, students will have developed a fluid, AI-augmented approach to research, equipping them to accelerate discovery, adapt to new research fields quickly, and tackle complex physical phenomena.
\end{itemize}

\subsection{Project: Majorana Modes in the Machine – Simulating Topological Phases with Quantum Circuits and AI}

In this project, your goal is to bridge the gap between abstract theoretical physics and the quantum software stack used to study many-body systems, 
utilizing LLMs as your active research collaborators. 
You will task AI with helping you simulate the topological phase transition of the 1D Kitaev chain – mapping fermions to qubits, 
preparing small-system ground states, and measuring observables that reveal the emergence of Majorana edge physics. 
Rather than relying on fragile access to real quantum hardware, you will build a controlled simulation environment on a classical computer, combining exact diagonalization, free-fermion methods, and quantum-circuit emulation.

Over 14 weeks, teams will progress through four key blocks:
\begin{enumerate}
\item \textbf{The Physics Bridge:} Building the Kitaev chain, finite-size spectra, and topological diagnostics.
\item \textbf{The Qubit Encoding:} Fermion-to-qubit mapping and ideal quantum-circuit simulation.
\item \textbf{Measuring Topology:} Designing circuit-based measurements for edge and string observables.
\item \textbf{The NISQ Reality Check:} Noise modeling, mitigation, and the robustness of topological signatures.
\end{enumerate}


